\section{Conclusion}

This study presents a systematic, reproducible evaluation of eight frontier large language models on stratified samples of USMLE Step 1 and Step 2 CK questions, with dedicated analysis of model performance from the perspective of international medical graduate students. Using a unified evaluation pipeline with standardized prompting, disk-cached API responses, and rigorous statistical methodology—including McNemar's tests with Holm-Bonferroni correction and Cohen's $h$ effect sizes—we have established a replicable framework for comparing LLM performance on high-stakes medical examination content.

Results to be reported after running the full evaluation pipeline will quantify the relative strengths and weaknesses of each model across medical subjects, difficulty levels, and exam steps, and will characterize the degree to which accuracy differs between questions containing US-centric clinical assumptions and those that do not.

Regardless of the specific numerical findings, the key practical takeaway from this work is clear: frontier LLMs have become capable enough to be genuinely useful study aids for USMLE preparation, but they are not uniformly reliable across all question types, and their reliability on content most relevant to IMG students remains uncertain. Medical students—particularly international graduates preparing for Step 2 CK—should treat AI-generated explanations as a starting point for understanding, not an endpoint, and should verify model outputs against authoritative clinical references.

For educators and program directors, this work suggests the value of developing formal guidelines for AI tool use in board preparation, analogous to the evidence-based guidance that exists for commercial question banks. For researchers, it highlights the need for larger-scale, pre-registered evaluations using validated items and diverse prompting strategies to generate the level of evidence that clinical guidance on AI study tools requires.
